\documentclass[11pt,a4paper]{article}
\usepackage[utf8]{inputenc}
\usepackage[T1]{fontenc}
\usepackage{amsmath,amssymb}
\usepackage{graphicx}
\usepackage{booktabs}
\usepackage{hyperref}
\usepackage{subcaption}
\usepackage{float}
\usepackage{xcolor}
\usepackage[margin=2.5cm]{geometry}

\title{Parameter Sweep Report: 2D Nonlinear Shallow Water\\
\large \texttt{smooth\_weight} $\times$ \texttt{base} Grid Search with Extended Rollout ($T_F = 2.0$)}
\author{}
\date{\today}

\begin{document}
\maketitle

\section{Motivation}

The previous experiment used \texttt{compare\_TF=1.0} (1-second rollout) with fixed hyperparameters
\texttt{smooth\_weight=0.05} and \texttt{base=32}.
Two changes were investigated simultaneously:
\begin{enumerate}
  \item \textbf{Longer rollout}: extend \texttt{compare\_TF} from 1.0 to 2.0\,s and increase output frames from 6 to 8, to better stress-test the NN surrogate's long-horizon prediction.
  \item \textbf{Hyperparameter search}: systematically vary the TV regularization weight \texttt{smooth\_weight} and CNN channel width \texttt{base} to find the configuration with the lowest L1 error over the 2-second rollout.
\end{enumerate}

\section{Experimental Setup}

\subsection{Model and training}

The surrogate is a residual CNN (3 input channels: $h - h_0$, $q_x$, $q_y$; 3 output channels) trained on
patches extracted from a Fourier-spectral shallow-water solver.
The loss function is
\begin{equation}
  \mathcal{L} = \underbrace{\frac{1}{N}\sum_i \lvert \hat{y}_i - y_i \rvert}_{\text{L1 data}}
  + \lambda \underbrace{\bigl(\lvert \nabla_x \hat{y}\rvert + \lvert \nabla_y \hat{y}\rvert\bigr)}_{\text{TV regularization}},
\end{equation}
where $\lambda = \texttt{smooth\_weight}$.
Training runs for 150 epochs with the learning rate schedule
$[(60,\,3\times10^{-4}),\,(110,\,10^{-4}),\,(140,\,3\times10^{-5}),\,(150,\,10^{-5})]$.

\subsection{Parameter grid}

Six combinations were evaluated (Table~\ref{tab:grid}).

\begin{table}[H]
  \centering
  \caption{Parameter grid: 6 combinations of \texttt{smooth\_weight} and \texttt{base}.}
  \label{tab:grid}
  \begin{tabular}{ccc}
    \toprule
    Combination & \texttt{smooth\_weight} $\lambda$ & \texttt{base} (CNN channels) \\
    \midrule
    1 & 0.0  & 32 \\
    2 & 0.01 & 32 \\
    3 & 0.05 & 32 \\
    4 & 0.1  & 32 \\
    5 & 0.0  & 48 \\
    6 & 0.05 & 48 \\
    \bottomrule
  \end{tabular}
\end{table}

All other settings are held fixed: grid $256\times256$, $n_{\mathrm{wd}}=32$, 6000 training samples,
batch size 128, warmup $T_{\mathrm{warmup}}=4.0$\,s, compare IC = \texttt{random}, seed 42.

\subsection{Evaluation metric}

After training, each model is rolled out for $T_F = 2.0$\,s alongside the spectral reference.
The comparison is sampled at 8 uniformly-spaced frames ($t_0,\ldots,t_7$).
The primary metric is the \textbf{mean L1 error} averaged over all frames and all three state components:
\begin{equation}
  \bar{\varepsilon}_{L1} = \frac{1}{N_t}\sum_{k=0}^{N_t-1}
  \frac{1}{|\Omega|}\sum_{(i,j)} \bigl\lvert U^{\mathrm{NN}}_{k,ij} - U^{\mathrm{ref}}_{k,ij}\bigr\rvert.
\end{equation}

\section{Results}

\subsection{Summary table}

Table~\ref{tab:results} lists all six combinations sorted by mean L1 error (ascending).

\begin{table}[H]
  \centering
  \caption{Sweep results sorted by mean L1 error (2-second rollout, 8 frames).}
  \label{tab:results}
  \begin{tabular}{ccccc}
    \toprule
    Rank & \texttt{smooth\_weight} & \texttt{base} & L1 mean & L1 max frame \\
    \midrule
    \textbf{1} & \textbf{0.05} & \textbf{48} & \textbf{0.014217} & \textbf{0.025887} \\
    2 & 0.1  & 32 & 0.019245 & 0.038809 \\
    3 & 0.0  & 32 & 0.020247 & 0.040786 \\
    4 & 0.01 & 32 & 0.024519 & 0.047638 \\
    5 & 0.05 & 32 & 0.031534 & 0.068902 \\
    6 & 0.0  & 48 & 0.035519 & 0.098977 \\
    \bottomrule
  \end{tabular}
\end{table}

\subsection{Key findings}

\begin{enumerate}
  \item \textbf{Best configuration: $\lambda=0.05$, \texttt{base}=48.}
    Combining moderate TV regularization with a wider CNN achieves the lowest error
    (L1 mean = 0.014217, max frame = 0.025887),
    a \textbf{30\% improvement} over the previous baseline ($\lambda=0.05$, \texttt{base}=32, L1 = 0.031534 on the same 2-second task).

  \item \textbf{TV regularization is essential for wider models.}
    Without regularization, \texttt{base}=48 is the \emph{worst} configuration (L1 = 0.035519, max = 0.098977),
    suggesting that the larger model memorizes short-horizon training patterns but diverges
    over longer rollouts. Adding $\lambda=0.05$ brings it to rank 1.

  \item \textbf{For \texttt{base}=32, less regularization is better.}
    Among the four base-32 runs, no regularization ($\lambda=0$) outperforms
    moderate regularization ($\lambda=0.05$, $\lambda=0.1$).
    This indicates the smaller model is already underfitting the spatial structure,
    and additional TV penalty further degrades its expressiveness.

  \item \textbf{Intermediate $\lambda$ values ($0.01$, $0.05$ at base=32) are suboptimal.}
    They neither fully regularize (as $\lambda=0.1$ does) nor allow full data-fitting
    (as $\lambda=0$ does), yielding middle-ground errors.
\end{enumerate}

\subsection{Runtime}

The best configuration (\texttt{base}=48) adds modest overhead over the base-32 model.
Table~\ref{tab:runtime} shows the final compare runtimes.

\begin{table}[H]
  \centering
  \caption{Rollout runtimes for the best configuration ($\lambda=0.05$, \texttt{base}=48, $T_F=2.0$\,s).}
  \label{tab:runtime}
  \begin{tabular}{lcc}
    \toprule
    Solver & Wall time (s) & Speedup \\
    \midrule
    Spectral (reference) & 71.076 & $1\times$ \\
    CNN surrogate (NN)   & 4.240  & $16.8\times$ \\
    \bottomrule
  \end{tabular}
\end{table}

\section{Comparison figures (best configuration: absolute TV, $\lambda=0.05$, \texttt{base}=48)}

Figures~\ref{fig:t0}--\ref{fig:t7} show the spectral reference vs.\ NN rollout
at 8 uniformly-spaced times over $[0,\,2.0]$\,s for state components $(h-h_0,\,q_x,\,q_y)$.

\begin{figure}[H]
  \centering
  \includegraphics[width=0.95\textwidth]{figures/wave_2d_nonlinear/sweep_abs_sw0.05_b48/t0.png}
  \caption{Frame 0: $t \approx 0.0$\,s.}
  \label{fig:t0}
\end{figure}

\begin{figure}[H]
  \centering
  \includegraphics[width=0.95\textwidth]{figures/wave_2d_nonlinear/sweep_abs_sw0.05_b48/t1.png}
  \caption{Frame 1: $t \approx 0.29$\,s.}
\end{figure}

\begin{figure}[H]
  \centering
  \includegraphics[width=0.95\textwidth]{figures/wave_2d_nonlinear/sweep_abs_sw0.05_b48/t2.png}
  \caption{Frame 2: $t \approx 0.57$\,s.}
\end{figure}

\begin{figure}[H]
  \centering
  \includegraphics[width=0.95\textwidth]{figures/wave_2d_nonlinear/sweep_abs_sw0.05_b48/t3.png}
  \caption{Frame 3: $t \approx 0.86$\,s.}
\end{figure}

\begin{figure}[H]
  \centering
  \includegraphics[width=0.95\textwidth]{figures/wave_2d_nonlinear/sweep_abs_sw0.05_b48/t4.png}
  \caption{Frame 4: $t \approx 1.14$\,s.}
\end{figure}

\begin{figure}[H]
  \centering
  \includegraphics[width=0.95\textwidth]{figures/wave_2d_nonlinear/sweep_abs_sw0.05_b48/t5.png}
  \caption{Frame 5: $t \approx 1.43$\,s.}
\end{figure}

\begin{figure}[H]
  \centering
  \includegraphics[width=0.95\textwidth]{figures/wave_2d_nonlinear/sweep_abs_sw0.05_b48/t6.png}
  \caption{Frame 6: $t \approx 1.71$\,s.}
\end{figure}

\begin{figure}[H]
  \centering
  \includegraphics[width=0.95\textwidth]{figures/wave_2d_nonlinear/sweep_abs_sw0.05_b48/t7.png}
  \caption{Frame 7: $t = 2.0$\,s (final frame).}
  \label{fig:t7}
\end{figure}

\section{Conclusion}

The hyperparameter sweep over 6 combinations identifies
\textbf{$\lambda = 0.05$, \texttt{base}=48} as the best configuration for the
2-second nonlinear shallow-water rollout, achieving mean L1 error of 0.014217
at a $16.8\times$ speedup over the spectral reference.
The interaction between model capacity (\texttt{base}) and TV regularization ($\lambda$)
is non-trivial: larger models require regularization to generalize over long horizons,
while smaller models are harmed by excessive TV penalty.

% ============================================================
\section{Per-channel Relative TV Regularization (Follow-up Sweep)}
\label{sec:relative_tv}
% ============================================================

\subsection{Motivation}

The three state channels $(h-h_0,\,q_x,\,q_y)$ differ in magnitude by roughly 50$\times$
(std$\,\approx 0.004$ for $h-h_0$, std$\,\approx 0.2$ for $q_x,q_y$).
The absolute TV penalty $\lambda(\lvert\nabla_x\hat{y}\rvert+\lvert\nabla_y\hat{y}\rvert)$
therefore penalizes $q_x,q_y$ far more than $h-h_0$.
A \emph{per-channel relative} formulation normalizes each channel's TV by its own predicted amplitude:
\begin{equation}
  \mathcal{L}_{TV}^{\text{rel}} = \lambda \sum_{c=0}^{2}
  \frac{\mathrm{TV}(\hat{y}_c)}{\mathbb{E}[\lvert\hat{y}_c\rvert] + \varepsilon},
  \qquad \varepsilon = 10^{-6},
\end{equation}
so that each channel receives an equal \emph{proportional} smoothing penalty.

\subsection{Experimental setup}

Model capacity is fixed at \texttt{base}=48 (best from the previous sweep).
Five values of $\lambda$ are evaluated: $\{0.01, 0.05, 0.1, 0.2, 0.5\}$.
All other settings (grid, epochs, LR schedule, seed, rollout length) are identical to
the absolute TV sweep in Section~\ref{sec:results}.
The sweep was stopped after 3 of 5 configurations completed (sw $= 0.2$ and $0.5$ not run)
because the trend was already conclusive.

\subsection{Results}

Table~\ref{tab:rel_results} shows the completed relative TV runs alongside all
absolute TV runs from Table~\ref{tab:results}.

\begin{table}[H]
  \centering
  \caption{Combined results: absolute TV (Sweep 1) and relative TV (Sweep 2, \texttt{base}=48).
           Rows sorted by L1 mean (ascending). Best configuration highlighted.}
  \label{tab:rel_results}
  \begin{tabular}{lcccc}
    \toprule
    TV mode & $\lambda$ & \texttt{base} & L1 mean & L1 max frame \\
    \midrule
    \textbf{absolute} & \textbf{0.05} & \textbf{48} & \textbf{0.014217} & \textbf{0.025887} \\
    absolute & 0.1  & 32 & 0.019245 & 0.038809 \\
    absolute & 0.0  & 32 & 0.020247 & 0.040786 \\
    absolute & 0.01 & 32 & 0.024519 & 0.047638 \\
    relative & 0.01 & 48 & 0.023784 & 0.042709 \\
    absolute & 0.05 & 32 & 0.031534 & 0.068902 \\
    absolute & 0.0  & 48 & 0.035519 & 0.098977 \\
    relative & 0.1  & 48 & 0.040901 & 0.070742 \\
    relative & 0.05 & 48 & 0.043298 & 0.081023 \\
    \bottomrule
  \end{tabular}
\end{table}

\subsection{Key findings}

\begin{enumerate}
  \item \textbf{Relative TV does not improve on absolute TV.}
    The best relative TV run ($\lambda=0.01$, L1 = 0.023784) is \textbf{67\% worse}
    than the absolute TV baseline ($\lambda=0.05$, L1 = 0.014217).

  \item \textbf{Relative TV degrades with increasing $\lambda$.}
    L1 rises from 0.023784 at $\lambda=0.01$ to 0.043298 at $\lambda=0.05$
    and 0.040901 at $\lambda=0.1$, indicating that larger relative penalties
    over-smooth all channels and prevent the CNN from learning fine spatial structure.

  \item \textbf{Root cause: amplitude normalization amplifies noise for $h-h_0$.}
    The $h-h_0$ channel has very small amplitude ($\approx 0.004$), so dividing its TV
    by $\mathbb{E}[\lvert\hat{y}_0\rvert]$ inflates its contribution to the loss,
    disrupting training rather than balancing the channels as intended.

  \item \textbf{Final choice: absolute TV, $\lambda=0.05$, \texttt{base}=48.}
    This configuration remains the best across both sweeps and is used in
    \texttt{config/wave\_2d\_nonlinear\_config.py}.
\end{enumerate}

\subsection{Comparison figures}

\subsubsection{Best relative TV run ($\lambda=0.01$, \texttt{base}=48, L1=0.023784)}

Figures~\ref{fig:rel_sw001_t0}--\ref{fig:rel_sw001_t7} show the full 8-frame rollout
for the best relative TV configuration alongside the spectral reference.
Errors are visibly larger than the absolute TV baseline (cf.\ Figures~\ref{fig:t0}--\ref{fig:t7}).

\begin{figure}[H]
  \centering
  \includegraphics[width=0.95\textwidth]{figures/wave_2d_nonlinear/sweep_rel_sw0.01/t0.png}
  \caption{Relative TV $\lambda=0.01$: frame 0 ($t\approx 0.0$\,s).}
  \label{fig:rel_sw001_t0}
\end{figure}
\begin{figure}[H]
  \centering
  \includegraphics[width=0.95\textwidth]{figures/wave_2d_nonlinear/sweep_rel_sw0.01/t1.png}
  \caption{Relative TV $\lambda=0.01$: frame 1 ($t\approx 0.29$\,s).}
\end{figure}
\begin{figure}[H]
  \centering
  \includegraphics[width=0.95\textwidth]{figures/wave_2d_nonlinear/sweep_rel_sw0.01/t2.png}
  \caption{Relative TV $\lambda=0.01$: frame 2 ($t\approx 0.57$\,s).}
\end{figure}
\begin{figure}[H]
  \centering
  \includegraphics[width=0.95\textwidth]{figures/wave_2d_nonlinear/sweep_rel_sw0.01/t3.png}
  \caption{Relative TV $\lambda=0.01$: frame 3 ($t\approx 0.86$\,s).}
\end{figure}
\begin{figure}[H]
  \centering
  \includegraphics[width=0.95\textwidth]{figures/wave_2d_nonlinear/sweep_rel_sw0.01/t4.png}
  \caption{Relative TV $\lambda=0.01$: frame 4 ($t\approx 1.14$\,s).}
\end{figure}
\begin{figure}[H]
  \centering
  \includegraphics[width=0.95\textwidth]{figures/wave_2d_nonlinear/sweep_rel_sw0.01/t5.png}
  \caption{Relative TV $\lambda=0.01$: frame 5 ($t\approx 1.43$\,s).}
\end{figure}
\begin{figure}[H]
  \centering
  \includegraphics[width=0.95\textwidth]{figures/wave_2d_nonlinear/sweep_rel_sw0.01/t6.png}
  \caption{Relative TV $\lambda=0.01$: frame 6 ($t\approx 1.71$\,s).}
\end{figure}
\begin{figure}[H]
  \centering
  \includegraphics[width=0.95\textwidth]{figures/wave_2d_nonlinear/sweep_rel_sw0.01/t7.png}
  \caption{Relative TV $\lambda=0.01$: frame 7 ($t=2.0$\,s, final).}
  \label{fig:rel_sw001_t7}
\end{figure}

\subsubsection{Final-frame comparison across relative TV configurations}

Figure~\ref{fig:rel_final} shows the final frame ($t=2.0$\,s) for all three completed
relative TV runs side by side, illustrating how larger $\lambda$ leads to over-smoothing
and larger rollout error.

\begin{figure}[H]
  \centering
  \begin{subfigure}[t]{0.32\textwidth}
    \includegraphics[width=\textwidth]{figures/wave_2d_nonlinear/sweep_rel_sw0.01/t7.png}
    \caption{$\lambda=0.01$, L1=0.023784}
  \end{subfigure}
  \hfill
  \begin{subfigure}[t]{0.32\textwidth}
    \includegraphics[width=\textwidth]{figures/wave_2d_nonlinear/sweep_rel_sw0.05/t7.png}
    \caption{$\lambda=0.05$, L1=0.043298}
  \end{subfigure}
  \hfill
  \begin{subfigure}[t]{0.32\textwidth}
    \includegraphics[width=\textwidth]{figures/wave_2d_nonlinear/sweep_rel_sw0.1/t7.png}
    \caption{$\lambda=0.1$, L1=0.040901}
  \end{subfigure}
  \caption{Final frame ($t=2.0$\,s) for three relative TV configurations (\texttt{base}=48).
           All are worse than the absolute TV baseline (L1=0.014217).}
  \label{fig:rel_final}
\end{figure}

\end{document}
